% ============================================================
%  SECTION 2 — THEORETICAL FOUNDATION
% ============================================================
\section{Theoretical Foundation}
\label{sec:theory}

% ------------------------------------------------------------
\subsection{Fractal Analysis and Dimension Estimation}
\label{sec:theory:fractal}

Fractal sets exhibit \textit{statistical self-similarity}, maintaining structural complexity across observation scales~\cite{Mandelbrot1982}. This property is characterized by power-law scaling, where a characteristic measure $M(\varepsilon)$ obtained at scale $\varepsilon$ satisfies

\begin{equation}
  M(\varepsilon) \;\sim\; C \cdot \varepsilon^{-D},
  \label{eq:powerlaw}
\end{equation}

\noindent with $C$ acting as a proportionality constant and $D$ representing the fractal dimension~\cite{Falconer2003}. Taking logarithms of both sides yields a linear relationship:

\begin{equation}
  \log M(\varepsilon) = \log C - D \log \varepsilon.
  \label{eq:loglog}
\end{equation}

\noindent $D$ is estimated as the negative slope of a least-squares regression over the pairs $\{(\log \varepsilon_i, \log M(\varepsilon_i))\}_{i=1}^{m}$ across $m$ discrete scales.

As the theoretical Hausdorff dimension~\cite{Hausdorff1919} is intractable on discrete grids, the Minkowski--Bouligand (box-counting) dimension $D_B$ serves as the standard approximation~\cite{Falconer2003}:

\begin{equation}
  D_B = \lim_{\varepsilon \to 0}
        \frac{\log N(\varepsilon)}{\log(1/\varepsilon)},
  \label{eq:boxdim}
\end{equation}

\noindent where $N(\varepsilon)$ is the minimum number of boxes of side $\varepsilon$ needed to cover the set. The regression's coefficient of determination ($R^{2}$) validates the presence of consistent power-law scaling, a prerequisite for interpreting $D_B$ as a meaningful structural descriptor~\cite{Lopes2011}.

% ------------------------------------------------------------
\subsection{Classical Box-Counting Method}
\label{sec:theory:bc}

The Classical Box-Counting (BC) method discretizes Eq.~\eqref{eq:boxdim} for binary images~\cite{Mandelbrot1982,Liebovitch1989}. Given an image $I$ of $N = W \times H$ pixels, the domain is partitioned into $\lceil W/\varepsilon \rceil \times \lceil H/\varepsilon \rceil$ disjoint cells of side $\varepsilon$. A cell is considered \textit{occupied} if it contains at least one foreground pixel:

\begin{equation}
  N(\varepsilon)
    = \sum_{i=0}^{\lceil W/\varepsilon \rceil - 1}
      \sum_{j=0}^{\lceil H/\varepsilon \rceil - 1}
      \mathbf{1}\!\left[
        \exists\,(x,y) \in B_{ij}(\varepsilon) : I(x,y) = 1
      \right].
  \label{eq:bc_count}
\end{equation}

Applying BC to grayscale medical images requires prior binarization, typically via Otsu's global thresholding~\cite{Otsu1979}. While binarization reduces the representational complexity, it irreversibly discards intensity gradients, making $D_B$ highly sensitive to the threshold parameter $t^{*}$.

To bypass exhaustive pixel scanning, optimized implementations pre-compute the \textit{integral image} $P[x][y] = \sum_{i \leq x, j \leq y} I[i][j]$ via the recurrence~\cite{Crow1984}:

\begin{equation}
  P[x][y]
    = I[x][y] + P[x{-}1][y] + P[x][y{-}1] - P[x{-}1][y{-}1].
  \label{eq:prefix}
\end{equation}

\noindent The foreground pixel count within any rectangular cell is then extracted in $\mathcal{O}(1)$ time:

\begin{equation}
  \sigma(x_1,y_1,x_2,y_2) =
    \begin{aligned}[t]
      &P[x_2][y_2] - P[x_1{-}1][y_2] \\
      &- P[x_2][y_1{-}1] + P[x_1{-}1][y_1{-}1].
    \end{aligned}
  \label{eq:rect_sum}
\end{equation}

% ------------------------------------------------------------
\subsection{Gliding Box Method}
\label{sec:theory:gb}

The Gliding Box (GB) method~\cite{Allain1991} evaluates a dense sliding window rather than a disjoint grid. A box of side $L$ is centered at every valid pixel position $(x, y)$, defining the \textit{box mass} as:

\begin{equation}
  m_{x,y}(L) = \sum_{i=0}^{L-1} \sum_{j=0}^{L-1} I[x+i][y+j].
  \label{eq:gb_mass}
\end{equation}

\noindent Across $T(L) = (W-L+1)(H-L+1)$ valid placements, the \textit{mass-frequency distribution} $n(m, L)$ dictates the normalized probability $P(m, L) = n(m, L) / T(L)$. $D_B$ is estimated by substituting the mean box mass into Eq.~\eqref{eq:loglog}~\cite{Plotnick1996}:

\begin{equation}
  \mu(L) = \sum_{m=0}^{L^{2}} m \, P(m, L).
  \label{eq:gb_mu1}
\end{equation}

GB also extracts \textit{lacunarity}, capturing the translational heterogeneity of the spatial mass distribution:

\begin{equation}
  \Lambda(L)
    = \frac{\mu^{(2)}(L) - [\mu(L)]^{2}}{[\mu(L)]^{2}},
  \label{eq:lacunarity}
\end{equation}

\noindent where $\mu^{(2)}(L) = \sum_{m} m^{2} P(m, L)$. Homogeneous textures yield $\Lambda \approx 0$, whereas clustered geometries produce large $\Lambda$~\cite{Allain1991,Tolle2008}. Integral images (Eq.~\eqref{eq:rect_sum}) optimize mass evaluations to $\mathcal{O}(1)$~\cite{Crow1984}, a critical reduction given the dense $\Theta(N)$ overlapping evaluations per scale.

% ------------------------------------------------------------
\subsection{Differential Box-Counting Method}
\label{sec:theory:dbc}

The Differential Box-Counting (DBC) method~\cite{Sarkar1994} avoids binarization by modeling the grayscale image $I \in \{0, \ldots, G{-}1\}$ as a three-dimensional topographic relief. For a cell of side $\varepsilon$, the intensity axis is partitioned into discrete bands of height $\delta = \lceil G / \lceil W/\varepsilon \rceil \rceil$. The number of boxes required to encompass the intensity gradient within cell $(i, j)$ is:

\begin{equation}
  n_{ij}(\varepsilon)
    = \left\lceil \frac{I_{\max}^{ij} - I_{\min}^{ij}}{\delta}
      \right\rceil + 1,
  \label{eq:dbc_count}
\end{equation}

\noindent where $I_{\max}^{ij}$ and $I_{\min}^{ij}$ denote local extrema. The total covering count $N(\varepsilon) = \sum_{i,j} n_{ij}(\varepsilon)$ determines $D_B$. While DBC preserves the structural fidelity of grayscale inputs, optimized execution relies on 2D sparse tables to reduce the local extrema search from $\Theta(\varepsilon^2)$ to $\mathcal{O}(1)$ queries.

% ------------------------------------------------------------
\subsection{Blanket Method}
\label{sec:theory:bm}

The Blanket Method (BM)~\cite{Peleg1984} computes $D_B$ by quantifying the apparent area growth of the intensity surface under sequential morphological operations. The upper ($u_\varepsilon$) and lower ($b_\varepsilon$) \textit{blankets} at dilation level $\varepsilon$ are recursively defined:

\begin{align}
  u_\varepsilon(x,y) &= \max\!\Bigl(u_{\varepsilon-1}(x,y)+1,\;
    \max_{(s,t)\in\mathcal{N}} u_{\varepsilon-1}(s,t)\Bigr),
  \label{eq:blanket_u} \\
  b_\varepsilon(x,y) &= \min\!\Bigl(b_{\varepsilon-1}(x,y)-1,\;
    \min_{(s,t)\in\mathcal{N}} b_{\varepsilon-1}(s,t)\Bigr),
  \label{eq:blanket_b}
\end{align}

\noindent initialized at $u_0 = b_0 = I$, with $\mathcal{N}$ representing the 4-connected neighborhood. The enclosed \textit{blanket volume} is:

\begin{equation}
  V(\varepsilon)
    = \sum_{x,y} \bigl[u_\varepsilon(x,y) - b_\varepsilon(x,y)\bigr].
  \label{eq:blanket_vol}
\end{equation}

\noindent The apparent surface area $A(\varepsilon) = \bigl[V(\varepsilon) - V(\varepsilon{-}1)\bigr] / 2$ substitutes $M(\varepsilon)$ in Eq.~\eqref{eq:loglog} to derive the fractal dimension. Operating strictly on intensity profiles, BM avoids thresholding artifacts. While a naive implementation incurs a neighborhood-dependent computational cost, separable monotonic deque filters~\cite{Lemire2006} optimize the 2D morphological extrema to $\Theta(N)$ total operations per scale.